\documentclass[12pt,a4paper]{article}
\usepackage[margin=2cm]{geometry}
\usepackage{amsmath,amssymb,mathtools,bm}
\usepackage{graphicx}
\usepackage{float}
\usepackage[T1]{fontenc}
\usepackage{lmodern}
\usepackage{xcolor}
\usepackage[colorlinks=true,linkcolor=blue!55!black]{hyperref}
\numberwithin{equation}{section}
\setlength{\parindent}{0pt}
\setlength{\parskip}{0.55em}
\linespread{1.2}
\newcommand{\M}{\mathcal M}
\newcommand{\ii}{\mathrm i}
\newcommand{\slsh}[1]{\not\!#1}
\title{Tree-Level $e^-\mu^-\to e^-\mu^-$ Scattering}
\author{Caiyi, Jiang}
\date{}
\begin{document}
\maketitle

\section{The process and the amplitude}
Consider the process:
\begin{equation}
 e^-(p_1)+\mu^-(p_2)\longrightarrow e^-(p'_1)+\mu^-(p'_2),
 \qquad p_1+p_2=p'_1+p'_2.
\end{equation}
According to the momentum conservation, the momentum transfer carried by the virtual photon is
\begin{equation}
 q=p'_1-p_1=p_2-p'_2.
\end{equation}
We use $g^{\mu\nu}=\operatorname{diag}(1,-1,-1,-1)$, $m_e=0$, and
$m_\mu=m$. The tree-level diagram is the $t$-channel photon exchange, so $q^2=t$.

\begin{figure}[htbp]
 \centering
 \includegraphics[width=0.6\textwidth]{C:/Users/LENOVO/OneDrive/Desktop/tranverse spin/step-0/figures/e_mu_t_channel_attached.png}
 \caption{Tree-level $t$-channel photon exchange for
 $e^-(p_1)+\mu^-(p_2)\to e^-(p'_1)+\mu^-(p'_2)$. }
 \label{fig:e-mu-t-channel}
\end{figure}

According to the QED Feynman rules, for specified external spin labels
$s_1,s_2,s'_1,s'_2$,
\begin{align}
 i\M_{s'_1s'_2;s_1s_2}
 &=\left[\bar u_e(p'_1,s'_1)(-\ii e\gamma^\mu)u_e(p_1,s_1)\right]
 \frac{-\ii g_{\mu\nu}}{q^2}
 \left[\bar u_\mu(p'_2,s'_2)(-\ii e\gamma^\nu)u_\mu(p_2,s_2)\right]\\
 &=\frac{\ii e^2}{q^2}
 \left[\bar u_e(p'_1,s'_1)\gamma^\mu u_e(p_1,s_1)\right]
 \left[\bar u_\mu(p'_2,s'_2)\gamma_\mu u_\mu(p_2,s_2)\right].
\end{align}
Removing the overall $\ii$ gives the amplitude
\begin{equation}
 \boxed{\M_{s'_1s'_2;s_1s_2}=
 \frac{e^2}{q^2}
 [\bar u_e(p'_1,s'_1)\gamma^\mu u_e(p_1,s_1)]
 [\bar u_\mu(p'_2,s'_2)\gamma_\mu u_\mu(p_2,s_2)] .}
 \label{eq:original-amplitude}
\end{equation}


\section{Unpolarized calculation using traces}
\label{sec:unpolarized}
We first calculate the unpolarized result as a reference. Sum over unobserved final spins and average over the initial spins of the unpolarized incoming particles. Completeness gives
\begin{equation}
 \sum_s u_e(p,s)\bar u_e(p,s)=\slsh p,\qquad
 \sum_s u_\mu(p,s)\bar u_\mu(p,s)=\slsh p+m.
\end{equation}
There are two initial spin states for each incoming fermion, hence the average
factor $1/4$. With $q^4\equiv(q^2)^2$,
\begin{align}
 \overline{|\M|^2}
 &\equiv\frac14\sum_{\text{all spins}}|\M|^2\\
 &=\frac{e^4}{4q^4}
 \operatorname{Tr}[\slsh p'_1\gamma^\mu\slsh p_1\gamma^\nu]
 \operatorname{Tr}[(\slsh p'_2+m)\gamma_\mu(\slsh p_2+m)\gamma_\nu]\\
 &=\frac{e^4}{4q^4}L_e^{\mu\nu}L^{(\mu)}_{\mu\nu}.
 \label{eq:trace-start}
\end{align}
The label $(\mu)$ on $L^{(\mu)}$ identifies the muon tensor; the repeated
Lorentz indices are contracted.

\subsection{Evaluate the muon trace}
Expand the two factors $(\slsh p'_2+m)$ and $(\slsh p_2+m)$:
\begin{align}
 L^{(\mu)}_{\mu\nu}
 &=\operatorname{Tr}[\slsh p'_2\gamma_\mu\slsh p_2\gamma_\nu]
 +m\operatorname{Tr}[\slsh p'_2\gamma_\mu\gamma_\nu]
 +m\operatorname{Tr}[\gamma_\mu\slsh p_2\gamma_\nu]
 +m^2\operatorname{Tr}[\gamma_\mu\gamma_\nu].
\end{align}
The middle two terms contain an odd number of gamma matrices, so their traces
vanish. Use
\begin{align}
 \operatorname{Tr}(\gamma^\alpha\gamma^\mu\gamma^\beta\gamma^\nu)
 &=4(g^{\alpha\mu}g^{\beta\nu}-g^{\alpha\beta}g^{\mu\nu}
       +g^{\alpha\nu}g^{\mu\beta}),\\
 \operatorname{Tr}(\gamma^\mu\gamma^\nu)&=4g^{\mu\nu}.
\end{align}
Writing $\slsh p'_2=p'_{2\alpha}\gamma^\alpha$ and
$\slsh p_2=p_{2\beta}\gamma^\beta$, the four-gamma term becomes
\begin{align}
 \operatorname{Tr}[\slsh p'_2\gamma^\mu\slsh p_2\gamma^\nu]
 &=4p'_{2\alpha}p_{2\beta}
 (g^{\alpha\mu}g^{\beta\nu}-g^{\alpha\beta}g^{\mu\nu}
       +g^{\alpha\nu}g^{\mu\beta})\\
 &=4(p_2'{}^\mu p_2^\nu+p_2'{}^\nu p_2^\mu
       -g^{\mu\nu}p'_2\!\cdot p_2).
\end{align}
The remaining two-gamma term is
$m^2\operatorname{Tr}(\gamma^\mu\gamma^\nu)=4m^2g^{\mu\nu}$.
Combining them gives
\begin{equation}
 L^{(\mu)\mu\nu}=4\left[p_2'{}^\mu p_2^\nu+p_2'{}^\nu p_2^\mu
 -g^{\mu\nu}(p'_2\!\cdot p_2-m^2)\right].
 \label{eq:muon-trace}
\end{equation}

\subsection{Evaluate the electron trace}
Since $m_e=0$, the electron trace is
\begin{align}
 L_e^{\mu\nu}
 &=\operatorname{Tr}[\slsh p'_1\gamma^\mu\slsh p_1\gamma^\nu]\\
 &=p'_{1\alpha}p_{1\beta}
   \operatorname{Tr}[\gamma^\alpha\gamma^\mu\gamma^\beta\gamma^\nu]\\
 &=4\left[p_1'{}^\mu p_1^\nu+p_1'{}^\nu p_1^\mu
       -g^{\mu\nu}(p'_1\!\cdot p_1)\right].
 \label{eq:electron-trace}
\end{align}

\subsection{Contract the traces}
For clarity, define $a=p'_1\!\cdot p_1$ and
$b=p'_2\!\cdot p_2$. Substituting Eqs.~\eqref{eq:muon-trace} and
\eqref{eq:electron-trace} gives
\begin{align}
 L_e^{\mu\nu}L^{(\mu)}_{\mu\nu}
 =16[&p_1'{}^\mu p_1^\nu+p_1'{}^\nu p_1^\mu-g^{\mu\nu}a]\\
 \times[&p'_{2\mu}p_{2\nu}+p'_{2\nu}p_{2\mu}
             -g_{\mu\nu}(b-m^2)].
\end{align}
Multiply each term in the first bracket by each term in the second. The four
momentum-by-momentum products give
\begin{equation}
 2(p'_1\!\cdot p'_2)(p_1\!\cdot p_2)
 +2(p'_1\!\cdot p_2)(p_1\!\cdot p'_2).
\end{equation}
The terms containing a metric give, in order,
\begin{align}
 &-2a(b-m^2)-2ab+4a(b-m^2)\\
 &=-2ab+2am^2-2ab+4ab-4am^2\\
 &=-2am^2.
\end{align}
Therefore
\begin{align}
 L_e^{\mu\nu}L^{(\mu)}_{\mu\nu}
 &=16\left[2(p'_1\!\cdot p'_2)(p_1\!\cdot p_2)
 +2(p'_1\!\cdot p_2)(p_1\!\cdot p'_2)-2m^2(p'_1\!\cdot p_1)\right]\\
 &=32\left[(p_1\!\cdot p'_2)(p'_1\!\cdot p_2)
 +(p_1\!\cdot p_2)(p'_1\!\cdot p'_2)-m^2(p_1\!\cdot p'_1)\right].
\end{align}
Insert this result into Eq.~\eqref{eq:trace-start}:
\begin{equation}
 \boxed{\overline{|\M|^2}=\frac{8e^4}{q^4}
 \left[(p_1\!\cdot p'_2)(p'_1\!\cdot p_2)
 +(p_1\!\cdot p_2)(p'_1\!\cdot p'_2)-m^2(p_1\!\cdot p'_1)\right].}
 \label{eq:invariant-result}
\end{equation}
This is the invariant result from the original calculation.

\subsection{Centre-of-momentum kinematics and final result}
\label{sec:cm-kinematics}
In the centre-of-momentum frame, the incoming electron travels along $+z$ and the incoming muon along $-z$. The outgoing electron is scattered through angle $\theta$ in the $xz$ plane.
\begin{figure}[H]
 \centering
 \includegraphics[width=0.6\textwidth]{C:/Users/LENOVO/OneDrive/Desktop/tranverse spin/step-0/figures/e_mu_cm_kinematics_attached.png}
 \caption{Centre-of-momentum geometry.}
 \label{fig:e-mu-cm-kinematics}
\end{figure}
The four-momenta are
\begin{align}
 p_1&=(k,0,0,k),&p_2&=(E,0,0,-k),\\
 p'_1&=(k,k\sin\theta,0,k\cos\theta),&
 p'_2&=(E,-k\sin\theta,0,-k\cos\theta),
\end{align}
with $E=\sqrt{k^2+m^2}$. Using
$p\!\cdot r=p^0r^0-\bm p\!\cdot\bm r$, calculate the scalar products:
\begin{align}
 p_1\!\cdot p_2
 &=kE-(0,0,k)\!\cdot(0,0,-k)=kE+k^2=k(E+k),\\
 p'_1\!\cdot p'_2
 &=kE-(k\sin\theta,0,k\cos\theta)\!\cdot
       (-k\sin\theta,0,-k\cos\theta)\\
 &=kE+k^2(\sin^2\theta+\cos^2\theta)=k(E+k),\\
 p_1\!\cdot p'_2
 &=kE-(0,0,k)\!\cdot(-k\sin\theta,0,-k\cos\theta)
 =k(E+k\cos\theta),\\
 p'_1\!\cdot p_2
 &=kE-(k\sin\theta,0,k\cos\theta)\!\cdot(0,0,-k)
 =k(E+k\cos\theta),\\
 p_1\!\cdot p'_1
 &=k^2-(0,0,k)\!\cdot(k\sin\theta,0,k\cos\theta)
 =k^2(1-\cos\theta).
\end{align}
The momentum transfer is
\begin{align}
 q&=p'_1-p_1\\
 &=(k,k\sin\theta,0,k\cos\theta)-(k,0,0,k)\\
 &=(0,k\sin\theta,0,k(\cos\theta-1)),\\
 q^2&=-k^2\sin^2\theta-k^2(\cos\theta-1)^2\\
 &=-k^2[\sin^2\theta+\cos^2\theta-2\cos\theta+1]\\
 &=-2k^2(1-\cos\theta),\\
 q^4&=(q^2)^2=4k^4(1-\cos\theta)^2.
 \label{eq:q-kinematics}
\end{align}
Substitute the scalar products into Eq.~\eqref{eq:invariant-result}:
\begin{align}
 \overline{|\M|^2}
 &=\frac{8e^4}{q^4}\left[
 k^2(E+k\cos\theta)^2+k^2(E+k)^2
 -m^2k^2(1-\cos\theta)\right]\\
 &=\frac{8e^4k^2}{4k^4(1-\cos\theta)^2}
 \left[(E+k\cos\theta)^2+(E+k)^2
 -m^2(1-\cos\theta)\right].
\end{align}
Thus the final spin-averaged result is
\begin{equation}
 \boxed{\overline{|\M|^2}=
 \frac{2e^4}{k^2(1-\cos\theta)^2}
 \left[(E+k)^2+(E+k\cos\theta)^2
       -m^2(1-\cos\theta)\right].}
 \label{eq:final-result}
\end{equation}


\section{Helicity amplitudes}
\label{sec:helicity}
We now calculate individual helicity amplitudes and check their spin average against the unpolarized result in Eq.~\eqref{eq:final-result}.

\subsection{Helicity spinors and centre-of-momentum kinematics}
We label each fermion by its helicity $\lambda=+$ or $-$, corresponding to
$h=+1$ or $h=-1$ in formulas that use $h$; the physical helicity eigenvalue
is $h/2$. For a momentum direction $(\vartheta,\varphi)$, define the
 two-component helicity spinors
\begin{align}
 \xi_+(p)&=\begin{pmatrix}\cos(\vartheta/2)\\
 e^{\ii\varphi}\sin(\vartheta/2)\end{pmatrix},&
 \xi_-(p)&=\begin{pmatrix}-e^{-\ii\varphi}\sin(\vartheta/2)\\
 \cos(\vartheta/2)\end{pmatrix}.
 \label{eq:xi}
\end{align}
They represent positive and negative helicity, respectively. Their angular
dependence specifies how the spinor components change as the momentum
direction changes.

In the Weyl representation, positive-energy fermion spinors are
\begin{align}
 u_+(p)&=\begin{pmatrix}\sqrt{E_f-|\bm p|}\,\xi_+\\
                       \sqrt{E_f+|\bm p|}\,\xi_+\end{pmatrix},&
 u_-(p)&=\begin{pmatrix}\sqrt{E_f+|\bm p|}\,\xi_-\\
                       \sqrt{E_f-|\bm p|}\,\xi_-\end{pmatrix}.
 \label{eq:helicity-spinors}
\end{align}
For the electron, $E_f=|\bm p|=k$, so one Weyl component vanishes for each
helicity. For the muon, $E_f=E=\sqrt{k^2+m^2}$ and $|\bm p|=k$, so both
components remain. Thus the massless-electron approximation and the nonzero
muon mass enter through the spinor components here.

We use the centre-of-momentum geometry introduced in
Sec.~\ref{sec:cm-kinematics} and shown in Fig.~\ref{fig:e-mu-cm-kinematics}.
The four-momenta, including their time components, are
\begin{align}
 p_1&=(k,0,0,k),&p_2&=(E,0,0,-k),\\
 p'_1&=(k,k\sin\theta,0,k\cos\theta),&
 p'_2&=(E,-k\sin\theta,0,-k\cos\theta).
\end{align}
We choose the polar and azimuthal angles to be
\begin{align}
 (\vartheta_{p_1},\varphi_{p_1})&=(0,0),&
 (\vartheta_{p_2},\varphi_{p_2})&=(\pi,\pi),\\
 (\vartheta_{p'_1},\varphi_{p'_1})&=(\theta,0),&
 (\vartheta_{p'_2},\varphi_{p'_2})&=(\pi-\theta,\pi).
 \label{eq:angles}
\end{align}
At a pole, the azimuth is a phase convention for the spinor rather than a
physical feature of the momentum direction. Keeping this convention fixed
makes the relative phases between spin states consistent.

\subsection{Helicity amplitude and calculation strategy}
Replace the spin labels in Eq.~\eqref{eq:original-amplitude} by
$s_1=\lambda_1$, $s_2=\lambda_2$, $s'_1=\lambda_3$, and
$s'_2=\lambda_4$. The amplitude with all four helicity labels explicit is
\begin{equation}
 \boxed{\M_{\lambda_3\lambda_4;\lambda_1\lambda_2}
 =\frac{e^2}{q^2}
 [\bar u_{e,\lambda_3}(p'_1)\gamma^\mu u_{e,\lambda_1}(p_1)]
 [\bar u_{\mu,\lambda_4}(p'_2)\gamma_\mu u_{\mu,\lambda_2}(p_2)] .}
 \label{eq:helicity-amplitude}
\end{equation}
For a fixed helicity assignment, evaluate it in this order:
\begin{enumerate}
 \item Choose the four helicities.
 \item Substitute the corresponding external spinors from
       Eq.~\eqref{eq:helicity-spinors}.
 \item Evaluate the electron and muon currents.
 \item Contract their Lorentz indices to obtain the amplitude.
 \item Take the absolute square.
\end{enumerate}
\textbf{Each individual helicity amplitude is calculated using explicit spinors.} The completeness relation $\sum_s u(p,s)\bar u(p,s)=\slsh p+m$ was used in Sec.~\ref{sec:unpolarized} to sum over unobserved spins and average over unpolarized initial spins. Keeping the helicity amplitudes
separate also allows coherent superpositions of different helicities to
interfere, as required for transverse-spin states.

\subsection{Electron current}
Set $c=\cos(\theta/2)$ and $s=\sin(\theta/2)$. The incoming electron moves
along $+z$, giving
\begin{equation}
 \xi_+(p_1)=\begin{pmatrix}1\\0\end{pmatrix},\qquad
 \xi_-(p_1)=\begin{pmatrix}0\\1\end{pmatrix}.
\end{equation}
The outgoing electron has direction angles $(\theta,0)$, so
\begin{equation}
 \xi_+(p'_1)=\begin{pmatrix}c\\s\end{pmatrix},\qquad
 \xi_-(p'_1)=\begin{pmatrix}-s\\c\end{pmatrix}.
\end{equation}
Their Hermitian conjugates are
\begin{equation}
 \xi_+^\dagger(p'_1)=(c,s),\qquad
 \xi_-^\dagger(p'_1)=(-s,c).
\end{equation}
The Pauli matrices needed for the spatial current components are
\begin{equation}
 \sigma^1=\begin{pmatrix}0&1\\1&0\end{pmatrix},\quad
 \sigma^2=\begin{pmatrix}0&-\ii\\\ii&0\end{pmatrix},\quad
 \sigma^3=\begin{pmatrix}1&0\\0&-1\end{pmatrix}.
\end{equation}
For each incoming and outgoing helicity, define the two-component products
\begin{equation}
 X^0=\xi_{h'}^\dagger(p'_1)\xi_h(p_1),\qquad
 X^i=\xi_{h'}^\dagger(p'_1)\sigma^i\xi_h(p_1)\quad(i=1,2,3).
\end{equation}
As an example, for the $++$ assignment,
\begin{align}
 X^0&=(c,s)\begin{pmatrix}1\\0\end{pmatrix}=c,&
 X^1&=(c,s)\begin{pmatrix}0\\1\end{pmatrix}=s,\\
 X^2&=(c,s)\begin{pmatrix}0\\\ii\end{pmatrix}=\ii s,&
 X^3&=(c,s)\begin{pmatrix}1\\0\end{pmatrix}=c.
\end{align}
\textbf{For a massless positive-helicity electron}, the nonzero Weyl component gives
$\bar u\gamma^\mu u=2k(X^0,X^1,X^2,X^3)$, hence
\begin{equation}
 J^\mu_e(p'_1+,p_1+)=(2kc,2ks,2\ii ks,2kc).
 \label{eq:electron-current-example}
\end{equation}

For clarity, here is a second explicit example, with outgoing helicity $-$
and incoming helicity $+$. In this case
$\xi_-^\dagger(p'_1)=(-s,c)$ and $\xi_+(p_1)=(1,0)^T$, so
\begin{align}
 X^0&=(-s,c)\begin{pmatrix}1\\0\end{pmatrix}=-s,&
 X^1&=(-s,c)\begin{pmatrix}0\\1\end{pmatrix}=c,\\
 X^2&=(-s,c)\begin{pmatrix}0\\\ii\end{pmatrix}=\ii c,&
 X^3&=(-s,c)\begin{pmatrix}1\\0\end{pmatrix}=-s.
\end{align}
Thus $(X^0,X^1,X^2,X^3)_{-+}=(-s,c,\ii c,-s)$. Repeating the same
multiplication for the other assignments gives
\begin{align}
 (X^0,X^1,X^2,X^3)_{++}&=(c,s,\ii s,c),&
 (X^0,X^1,X^2,X^3)_{--}&=(c,-s,\ii s,-c),\\
 (X^0,X^1,X^2,X^3)_{+-}&=(s,c,-\ii c,-s),&
 (X^0,X^1,X^2,X^3)_{-+}&=(-s,c,\ii c,-s).
 \label{eq:electron-products}
\end{align}
The first sign labels the outgoing helicity and the second the incoming
helicity.

\paragraph{From the Pauli products to the full electron currents.}

Define $J^\mu_{e,h'h}=\bar u_{h'}(p'_1)\gamma^\mu u_h(p_1)$, with the
outgoing helicity first. In the massless limit, Eq.~\eqref{eq:helicity-spinors}
becomes
\begin{equation}
 u_+(p)=\sqrt{2k}\begin{pmatrix}0_2\\\xi_+(p)\end{pmatrix},\qquad
 u_-(p)=\sqrt{2k}\begin{pmatrix}\xi_-(p)\\0_2\end{pmatrix},
 \qquad 0_2=\begin{pmatrix}0\\0\end{pmatrix}.
\end{equation}
The matrices entering the current are block diagonal:
\begin{equation}
 \gamma^0\gamma^0=\begin{pmatrix}I_2&0\\0&I_2\end{pmatrix},\qquad
 \gamma^0\gamma^i=\begin{pmatrix}-\sigma^i&0\\0&\sigma^i\end{pmatrix},
 \qquad i=1,2,3.
\end{equation}
Since $J^\mu=u'^\dagger\gamma^0\gamma^\mu u$, the lower Weyl block gives
$J^0_{e,++}=2kX^0_{++}$ and $J^i_{e,++}=2kX^i_{++}$. The upper block gives
$J^0_{e,--}=2kX^0_{--}$ and $J^i_{e,--}=-2kX^i_{--}$: only the spatial
components acquire a minus sign. Substitution therefore gives
\begin{equation}
\boxed{\begin{aligned}
 J^\mu_{e,++}&=2k(X^0_{++},X^1_{++},X^2_{++},X^3_{++})
 \\[-0.2em]
 &=(2kc,2ks,2\ii ks,2kc),\\[0.6em]
 J^\mu_{e,--}&=2k(X^0_{--},-X^1_{--},-X^2_{--},-X^3_{--})
 \\[-0.2em]
 &=(2kc,2ks,-2\ii ks,2kc).
\end{aligned}}
 \label{eq:electron-currents}
\end{equation}
For a helicity-flip transition, the incoming and outgoing spinors occupy
opposite Weyl blocks. Every $\gamma^0\gamma^\mu$ above is block diagonal,
so it cannot connect those blocks. Explicitly, writing
$D_i=\operatorname{diag}(-\sigma^i,\sigma^i)$,
\begin{align}
 J^0_{e,+-}&=2k(0_2^\dagger,\xi_+'{}^\dagger)
 \begin{pmatrix}\xi_-\\0_2\end{pmatrix}=0,&
 J^i_{e,+-}&=2k(0_2^\dagger,\xi_+'{}^\dagger)D_i
 \begin{pmatrix}\xi_-\\0_2\end{pmatrix}=0,\\
 J^0_{e,-+}&=2k(\xi_-'{}^\dagger,0_2^\dagger)
 \begin{pmatrix}0_2\\\xi_+\end{pmatrix}=0,&
 J^i_{e,-+}&=2k(\xi_-'{}^\dagger,0_2^\dagger)D_i
 \begin{pmatrix}0_2\\\xi_+\end{pmatrix}=0.
\end{align}
\begin{equation}
      \boxed{J^\mu_{e,+-}=0,\qquad J^\mu_{e,-+}=0.}
\end{equation}
Here $\xi'_h=\xi_h(p'_1)$ and $\xi_h=\xi_h(p_1)$. This is the explicit
spinor form of helicity conservation at the massless electron vector vertex.
The Pauli products $X^\mu_{+-}$ and $X^\mu_{-+}$ can be nonzero, while the
complete currents vanish because of their Weyl-block structure.

\subsection{Muon current}
For the muon, both Weyl components contribute. 

Write
$u_h=(a_h\xi_h,b_h\xi_h)^T$, where
$a_h=\sqrt{E-hk}$ and $b_h=\sqrt{E+hk}$ for $h=+1,-1$. Multiplying the
upper and lower two-component blocks gives
\begin{align}
 \bar u_{h'}(p'_2)\gamma^0u_h(p_2)
 &=(b_{h'}b_h+a_{h'}a_h)\,\xi_{h'}^\dagger(p'_2)\xi_h(p_2),\\
 \bar u_{h'}(p'_2)\gamma^iu_h(p_2)
 &=(b_{h'}b_h-a_{h'}a_h)\,\xi_{h'}^\dagger(p'_2)\sigma^i\xi_h(p_2),
 \qquad i=1,2,3.
\end{align}
With $a_\pm=\sqrt{E\mp k}$ and $b_\pm=\sqrt{E\pm k}$, the coefficients
multiplying the time and spatial two-component products are
\begin{align}
 A_{++}=A_{--}&=b_+^2+a_+^2=2E,\\
 B_{++}&=b_+^2-a_+^2=2k,& B_{--}&=b_-^2-a_-^2=-2k,\\
 A_{+-}=A_{-+}&=b_+b_-+a_+a_-=2m,&
 B_{+-}&=B_{-+}=b_+b_--a_+a_-=0.
\end{align}
The resulting four muon currents are
\begin{equation}
\boxed{\begin{aligned}
 J^\mu_{\mu,++}&=(2Ec,-2ks,2\ii ks,-2kc),&
 J^\mu_{\mu,--}&=(2Ec,-2ks,-2\ii ks,-2kc),\\
 J^\mu_{\mu,+-}&=(2ms,0,0,0),&
 J^\mu_{\mu,-+}&=(-2ms,0,0,0).
\end{aligned}}
 \label{eq:muon-currents}
\end{equation}
Here the first helicity is outgoing and the second is incoming.
In the present centre-of-momentum kinematics, \textbf{the muon helicity-flip
currents have only a time component and are proportional to the mass $m$}.
They therefore vanish in the limit $m\to0$. 
\begin{itemize}
      \item The $++$ current at $\theta=0$ reduces to $(2E,0,0,-2k)$, as expected for a muon moving along $-z$.
\end{itemize}


\subsection{The helicity-amplitude table for this process}
The amplitude is the product of the electron and muon currents, 
\begin{equation}
      \M_{\lambda_3\lambda_4;\lambda_1\lambda_2}
      =\frac{e^2}{q^2}J_e\!\cdot J_\mu
\end{equation}
For example, for initial helicities $(\lambda_1,\lambda_2)=(+,+)$ and final
helicities $(\lambda_3,\lambda_4)=(+,+)$,
\begin{align}
 J_e\!\cdot J_\mu
 &=J_e^0J_\mu^0-J_e^1J_\mu^1-J_e^2J_\mu^2-J_e^3J_\mu^3\\
 &=(2kc)(2Ec)-(2ks)(-2ks)-(2\ii ks)(2\ii ks)
       -(2kc)(-2kc)\\
 &=4kEc^2+4k^2s^2+4k^2s^2+4k^2c^2\\
 &=4k\left[(E+k)c^2+2ks^2\right].
\end{align}
Define
\begin{equation}
 F=(E+k)c^2+2ks^2,\qquad G=(E+k)c^2,\qquad H=m\,s\,c.
\end{equation}
The same four multiplications for each allowed helicity choice give the table
below. Rows are the incoming pair $(\lambda_1,\lambda_2)$ and columns are
the outgoing pair $(\lambda_3,\lambda_4)$. The entries are $q^2\M/(4ke^2)$; multiply each by $4ke^2/q^2$ to recover the corresponding amplitude. 
\begin{equation}
\renewcommand{\arraystretch}{1.25}
\begin{array}{c|cccc}
 (\lambda_1,\lambda_2)\backslash(\lambda_3,\lambda_4)
 &++&+-&-+&--\\ \hline
 ++&F&-H&0&0\\
 +-&H&G&0&0\\
 -+&0&0&G&-H\\
 --&0&0&H&F
\end{array}
 \label{eq:helicity-table}
\end{equation}
The table contains
\begin{itemize}
 \item two entries equal to $F$;
 \item two entries equal to $G$;
 \item four entries with magnitude $H$;
 \item zeros in all remaining positions.
\end{itemize}
\textbf{Different final helicity states are orthogonal. }When the final spins are
unobserved, we therefore \textbf{sum their probabilities}. For unpolarized
incoming particles, we also average over the initial helicity states.
The spin average is defined by
\begin{equation}
 \overline{|\M|^2}_{\rm helicity}
 =\frac14\sum_{\lambda_1,\lambda_2}\sum_{\lambda_3,\lambda_4}
 \left|\M_{\lambda_3\lambda_4;\lambda_1\lambda_2}\right|^2.
\end{equation}
Every amplitude in the table has the common factor $4ke^2/q^2$.
For each fixed incoming pair, sum the squared magnitudes of the four
entries across its row. Since $F$, $G$, and $H$ are real, these row sums are
\begin{align}
 (\lambda_1,\lambda_2)=(+,+):\quad
 &F^2+(-H)^2+0+0=F^2+H^2,\\
 (\lambda_1,\lambda_2)=(+,-):\quad
 &H^2+G^2+0+0=H^2+G^2,\\
 (\lambda_1,\lambda_2)=(-,+):\quad
 &0+0+G^2+(-H)^2=G^2+H^2,\\
 (\lambda_1,\lambda_2)=(-,-):\quad
 &0+0+H^2+F^2=H^2+F^2.
\end{align}
\begin{align}
 \overline{|\M|^2}_{\rm helicity}
 &=\frac14\left(\frac{4ke^2}{q^2}\right)^2
       (2F^2+2G^2+4H^2)\\
 &=\frac{8k^2e^4}{q^4}(F^2+G^2+2H^2).
\end{align}
The factor $1/4$ arises because the incoming electron and muon each have
two spin states, giving four equally weighted initial helicity pairs.

Since $c^2=1-s^2$ and $\cos\theta=1-2s^2$, direct expansion gives
\begin{align}
 F^2+G^2+2H^2
 &=(E+k-2ks^2)^2+(E+k)^2-2m^2s^2\\
 &=(E+k\cos\theta)^2+(E+k)^2
       -m^2(1-\cos\theta).
\end{align}
Substituting this expression into the helicity spin average and using
$q^4=4k^4(1-\cos\theta)^2$ from Eq.~\eqref{eq:q-kinematics} gives the final result
\begin{equation}
 \boxed{\begin{aligned}
 \overline{|\M|^2}_{\rm helicity}
 &=\frac{2e^4}{k^2(1-\cos\theta)^2}
 &\quad\times\left[(E+k)^2+(E+k\cos\theta)^2
             -m^2(1-\cos\theta)\right].
 \end{aligned}}
 \label{eq:helicity-final-result}
\end{equation}
This agrees with the trace result in Eq.~\eqref{eq:final-result}, checking
the unpolarized normalization and the sum of squared helicity amplitudes.


\section{Transverse-spin states and interference}
\label{sec:transverse}
We now use the helicity amplitudes to describe a muon whose spin points
perpendicular to its momentum. The key rule is to combine amplitudes
first and take the absolute square afterwards.

\subsection{A transverse state contains both helicities}
A general normalized spin state is
\begin{equation}
 |\psi\rangle=\alpha|+\rangle+\beta|-\rangle,\qquad
 |\alpha|^2+|\beta|^2=1.
 \label{eq:helicity-superposition}
\end{equation}
Here $|\alpha|^2$ and $|\beta|^2$ are the two helicity probabilities.
The corresponding external spinor is $u_\psi=\alpha u_++\beta u_-$.
For a pure transverse state, the two probabilities are equal. Their
relative phase determines the transverse direction.

For our incoming muon, the chosen phases give
$\xi_+(p_2)=(0,-1)^T$ and $\xi_-(p_2)=(1,0)^T$.
Their difference is proportional to $(1,1)^T$, the spin state along $+x$.
Thus the physical $x$ states in this convention are
\begin{equation}
 u_{\mu,+x}=\frac{u_+-u_-}{\sqrt2},\qquad
 u_{\mu,-x}=\frac{u_++u_-}{\sqrt2}.
 \label{eq:x-spinors}
\end{equation}
All spinors here have momentum $p_2$. The minus sign for $+x$ follows from
this phase convention.

\subsection{Adding amplitudes and finding interference}
Fix the electron's incoming helicity and the final helicities. Let $A_+$
and $A_-$ be the amplitudes for incoming muon helicities $+$ and $-$.
Linearity gives
\begin{equation}
 \M_\psi=\alpha A_++\beta A_-.
\end{equation}
For the $+x$ state, this becomes
\begin{equation}
 \M_{+x}=\frac{A_+-A_-}{\sqrt2}.
 \label{eq:x-amplitude}
\end{equation}
For example, with incoming electron helicity $+$ and final state $++$,
the table gives $A_+=CF$ and $A_-=CH$, where $C=4ke^2/q^2$. Therefore,
\begin{equation}
 \M_{++;+,+x}=\frac{C}{\sqrt2}(F-H).
\end{equation}
Taking the absolute square gives
\begin{equation}
 |\M_{+x}|^2=\frac12\left(
 |A_+|^2+|A_-|^2-2\operatorname{Re}[A_+A_-^*]\right).
 \label{eq:interference}
\end{equation}
The last term is interference between two helicities leading to the same
final state. An unpolarized muon gives
$(|A_+|^2+|A_-|^2)/2$, without this interference term.

When final spins are unobserved, sum their probabilities. For incoming
electron helicities $+$ and $-$, respectively, the table gives
\begin{align}
 P_+&=\frac{C^2}{2}\left[(F-H)^2+(G+H)^2\right],\\
 P_-&=\frac{C^2}{2}\left[(G-H)^2+(F+H)^2\right].
\end{align}
For an unpolarized electron, average these two expressions:
\begin{equation}
 \frac{P_++P_-}{2}
 =\frac{C^2}{2}(F^2+G^2+2H^2)
 =\overline{|\M|^2}_{\rm unpolarized}.
\end{equation}
The interference contributions cancel. In our tree-level calculation,
polarizing only the muon transversely therefore leaves this spin-summed
result unchanged when the electron is unpolarized.

\subsection{Where the density matrix comes from}
A density matrix records spin probabilities and their interference.
In the ordered helicity basis $(|+\rangle,|-\rangle)$, the state is the
coefficient column $(\alpha,\beta)^T$. For a pure state, define
\begin{equation}
 \rho=|\psi\rangle\langle\psi|
 =\begin{pmatrix}\alpha\\\beta\end{pmatrix}(\alpha^*,\beta^*)
 =\begin{pmatrix}
 |\alpha|^2&\alpha\beta^*\\
 \alpha^*\beta&|\beta|^2
 \end{pmatrix}.
 \label{eq:pure-spin-density}
\end{equation}
The diagonal entries are the two helicity probabilities. The off-diagonal
entries contain the relative phase.

To write these entries using real numbers, define
\begin{equation}
 s_z=|\alpha|^2-|\beta|^2,\qquad
 s_x=2\operatorname{Re}(\alpha^*\beta),\qquad
 s_y=2\operatorname{Im}(\alpha^*\beta).
\end{equation}
Normalization then gives
\begin{equation}
 |\alpha|^2=\frac{1+s_z}{2},\qquad
 |\beta|^2=\frac{1-s_z}{2},\qquad
 \alpha\beta^*=\frac{s_x-\ii s_y}{2}.
\end{equation}
Substitution gives the density matrix
\begin{equation}
 \boxed{\rho=\frac12\begin{pmatrix}
 1+s_z&s_x-\ii s_y\\
 s_x+\ii s_y&1-s_z
 \end{pmatrix}.}
 \label{eq:spin-density-matrix}
\end{equation}
The parameters $s_x,s_y,s_z$ describe polarization in the helicity basis;
they are distinct from $s=\sin(\theta/2)$. For our phase convention, the
global $+x$ state has $(\alpha,\beta)=(1,-1)/\sqrt2$, so $s_x=-1$.
For comparison,
\begin{equation}
 \rho_{+x}=\frac12\begin{pmatrix}1&-1\\-1&1\end{pmatrix},\qquad
 \rho_{\rm unpolarized}=\frac12\begin{pmatrix}1&0\\0&1\end{pmatrix}.
\end{equation}
Both have equal helicity probabilities. Only the transverse pure state has
off-diagonal entries.

For a beam containing different spin states, average their density matrices
with their probabilities. The same matrix form applies. A fully polarized
state has $|\bm s|=1$, a partially polarized beam has $0<|\bm s|<1$, and an
unpolarized beam has $\bm s=0$.

For fixed electron helicity and final state, the squared amplitude obtained
with this matrix is
\begin{equation}
 \begin{aligned}
 \mathcal P_f(\rho)
 &=\rho_{++}|A_+|^2+\rho_{--}|A_-|^2\\
 &\quad+\rho_{+-}A_+A_-^*+\rho_{-+}A_-A_+^*.
 \end{aligned}
\end{equation}
The first two terms weight the helicity probabilities. The last two terms
give interference. Inserting $\rho_{+x}$ reproduces
$|A_+-A_-|^2/2$; inserting $\rho_{\rm unpolarized}$ gives
$(|A_+|^2+|A_-|^2)/2$.

\end{document}
